\begin{thebibliography}{99}
\addcontentsline{toc}{chapter}{Bibliographie}

\bibitem{alten} ALTEN Group, \emph{Site institutionnel et résultats annuels},
\url{https://www.alten.com} --- consulté en 2026.

\bibitem{altenlabs} ALTEN Group, \emph{Les Labs ALTEN : un temps d'avance sur
l'innovation}, \url{https://www.alten.com/fr/les-labs-alten-un-temps-davance-sur-linnovation/}.

\bibitem{sae} SAE International, \emph{J3016 : Taxonomy and Definitions for
Terms Related to Driving Automation Systems}, 2021.

\bibitem{ros2} Open Robotics, \emph{ROS 2 Documentation (Humble)},
\url{https://docs.ros.org}.

\bibitem{etsi-cam} ETSI, \emph{EN 302 637-2 : Intelligent Transport Systems ;
Cooperative Awareness Basic Service (CAM)}.

\bibitem{etsi-denm} ETSI, \emph{EN 302 637-3 : Intelligent Transport Systems ;
Decentralized Environmental Notification Basic Service (DENM)}.

\bibitem{vanetza} R. Riebl et al., \emph{Vanetza : an open-source
implementation of the ETSI ITS protocol stack},
\url{https://www.vanetza.org}.

\bibitem{zenoh} Eclipse Foundation, \emph{Zenoh : Zero Overhead Network
Protocol}, \url{https://zenoh.io}.

\bibitem{f1tenth} M. O'Kelly et al., \emph{F1TENTH : An Open-source Evaluation
Environment for Continuous Control and Reinforcement Learning}, NeurIPS 2019.

\bibitem{yolo} J. Redmon et al., \emph{You Only Look Once : Unified, Real-Time
Object Detection}, CVPR 2016.

\bibitem{nav2} S. Macenski et al., \emph{The Marathon 2 : A Navigation System},
IROS 2020.

% TODO : compléter au fil de la rédaction (3GPP 5G, slam_toolbox, articles
% téléopération...) — la grille INSA valorise une biblio réelle et citée
% dans le texte via \cite{...}.
\end{thebibliography}
